\documentclass[UTF8,a4paper,11pt]{ctexart}

\usepackage{amsmath,amssymb,mathtools}
\usepackage{geometry}
\usepackage{fancyhdr}
\usepackage{enumitem}
\usepackage{hyperref}

\geometry{left=2.5cm,right=2.5cm,top=2.5cm,bottom=2.5cm}
\setlength{\parindent}{2em}
\linespread{1.2}
\setlength{\headheight}{14pt}

\pagestyle{fancy}
\fancyhf{}
\fancyfoot[C]{\thepage}
\fancyhead[L]{数值分析\quad 第十一次理论作业\quad 学号：241098117，姓名：张城宗}
\fancyhead[R]{第五章}

\newcommand{\dd}{\,\mathrm{d}}
\newcommand{\ee}{\mathrm{e}}
\newcommand{\QED}{\hfill$\square$}

\begin{document}

\begin{center}
  {\LARGE\bfseries 数值分析\quad 第十一次理论作业}\\[6pt]
  {\large 第五章第 32、34(2)、36(2)(4) 题}\\[4pt]
  {\normalsize 学号：241098117\quad 姓名：张城宗}
\end{center}

\bigskip

%=====================================================
\noindent\textbf{第五章\quad 第 32 题}\quad
证明 $m$ 阶齐次线性差分方程 (5.5.7) 的线性无关解的最大数目是 $m$.

\medskip
\noindent\textbf{证明}\quad
$m$ 阶齐次线性差分方程为
\[
  a_0(n)y(n)+a_1(n)y(n+1)+\cdots+a_m(n)y(n+m)=0,
  \tag{5.5.7}
\]
其中 $a_0(n)\neq 0$，$a_m(n)\neq 0$. 据习题 31（性质 5.5.2），方程 (5.5.7) 至少有 $m$ 个线性无关解，故线性无关解的最大数目至少为 $m$. 下面证明至多为 $m$.

\textbf{至多 $m$ 个线性无关解.}\quad
设 $y_1(n),y_2(n),\ldots,y_{m+1}(n)$ 是方程 (5.5.7) 的任意 $m+1$ 个解. 考虑关于未知常数 $C_1,C_2,\ldots,C_{m+1}$ 的齐次线性方程组：
\[
  \sum_{j=1}^{m+1} C_j\, y_j(i)=0,\qquad i=0,1,\ldots,m-1.
\]
这是含 $m$ 个方程、$m+1$ 个未知量的齐次线性方程组，由线性方程组理论知必存在非零解，即存在
$(C_1,C_2,\ldots,C_{m+1})\neq\mathbf{0}$.

令
\[
  g(n)=C_1 y_1(n)+C_2 y_2(n)+\cdots+C_{m+1}y_{m+1}(n).
\]
由方程 (5.5.7) 的线性性，$g(n)$ 也满足 (5.5.7)，且由上面方程组的解的性质有
\[
  g(0)=g(1)=\cdots=g(m-1)=0.
\]
由于 $a_m(n)\neq 0$，方程 (5.5.7) 可改写为
\[
  y(n+m)=-\frac{1}{a_m(n)}\bigl[a_0(n)y(n)+\cdots+a_{m-1}(n)y(n+m-1)\bigr],
\]
故解由 $m$ 个连续初始值唯一确定. 既然 $g$ 在 $0,1,\ldots,m-1$ 处均为零，由解的唯一性得
\[
  g(n)=0 \quad \text{对所有 } n\in N.
\]
这说明 $y_1(n),y_2(n),\ldots,y_{m+1}(n)$ 线性相关. 因为 $m+1$ 个解任意取均线性相关，所以线性无关解的数目至多为 $m$.

综合两方面，$m$ 阶齐次线性差分方程 (5.5.7) 的线性无关解的最大数目恰好是 $m$. \QED

\bigskip\bigskip

%=====================================================
\noindent\textbf{第五章\quad 第 34 题 (2)}\quad
求差分方程
\[
  y(n+2)+2y(n+1)+2y(n)=2^n
\]
的通解.

\medskip
\noindent\textbf{解}

\textbf{第一步：求对应齐次方程的通解.}\quad
对应的齐次差分方程为
\[
  y(n+2)+2y(n+1)+2y(n)=0.
\]
其特征方程为
\[
  z^2+2z+2=0,
\]
解得
\[
  z=\frac{-2\pm\sqrt{4-8}}{2}=-1\pm i.
\]
写成极坐标形式：$|-1+i|=\sqrt{2}$，$\arg(-1+i)=\dfrac{3\pi}{4}$，故特征根为
\[
  z=\sqrt{2}\,\ee^{\pm i\frac{3\pi}{4}}=\sqrt{2}\!\left(\cos\frac{3\pi}{4}\pm i\sin\frac{3\pi}{4}\right).
\]
由公式 (5.5.18)，齐次方程的通解为
\[
  y_h(n)=(\sqrt{2})^n\!\left(C_1\cos\frac{3n\pi}{4}+C_2\sin\frac{3n\pi}{4}\right),
\]
其中 $C_1,C_2$ 为任意常数.

\textbf{第二步：求特解.}\quad
右端 $2^n$ 对应 $z_0=2$，而 $z_0=2$ 不是特征方程 $z^2+2z+2=0$ 的根. 设特解形式为
\[
  y^*(n)=K\cdot 2^n.
\]
代入方程：
\[
  K\cdot 2^{n+2}+2K\cdot 2^{n+1}+2K\cdot 2^n
  =\bigl(4+4+2\bigr)K\cdot 2^n
  =10K\cdot 2^n=2^n,
\]
故 $K=\dfrac{1}{10}$，即
\[
  y^*(n)=\frac{1}{10}\cdot 2^n.
\]

\textbf{第三步：写出通解.}\quad
由性质 5.5.4，非齐次方程的通解为齐次通解与特解之和：
\[
  \boxed{
  y(n)=(\sqrt{2})^n\!\left(C_1\cos\frac{3n\pi}{4}+C_2\sin\frac{3n\pi}{4}\right)+\frac{2^n}{10},
  }
\]
其中 $C_1,C_2$ 为任意常数.

\bigskip\bigskip

%=====================================================
\noindent\textbf{第五章\quad 第 36 题 (2)}\quad
判断解初值问题 $y'=f(t,y)$，$y(t_0)=y_0$ 的多步法
\[
  y_{n+2}+y_{n+1}-2y_n=h(2f_{n+1}+f_n)
\]
是否收敛，并说明原因.

\medskip
\noindent\textbf{解}\quad
将方法写成标准线性 $k$ 步法 (5.1.10) 的形式，这里 $k=2$，系数为
\[
  \alpha_2=1,\quad\alpha_1=1,\quad\alpha_0=-2;\qquad
  \beta_2=0,\quad\beta_1=2,\quad\beta_0=1.
\]
对应的特征多项式为
\[
  \rho(\lambda)=\lambda^2+\lambda-2=(\lambda-1)(\lambda+2),\qquad
  \sigma(\lambda)=2\lambda+1.
\]

\textbf{相容性检验.}\quad
\[
  \rho(1)=1+1-2=0\;\checkmark,\qquad
  \rho'(\lambda)=2\lambda+1,\quad\rho'(1)=3,\quad\sigma(1)=2+1=3\;\checkmark.
\]
相容条件 (5.6.8) 满足，故方法是\textbf{相容的}.

\textbf{零稳定性检验.}\quad
$\rho(\lambda)=(\lambda-1)(\lambda+2)$ 的根为 $\lambda_1=1$ 和 $\lambda_2=-2$.
由于 $|\lambda_2|=2>1$，存在模大于 $1$ 的根，\textbf{特征根条件不满足}.

由定理 5.6.5，方法\textbf{不稳定（不满足零稳定性）}.

\textbf{结论.}\quad
方法不满足零稳定性，因此\textbf{不收敛}.

\bigskip\bigskip

%=====================================================
\noindent\textbf{第五章\quad 第 36 题 (4)}\quad
判断解初值问题 $y'=f(t,y)$，$y(t_0)=y_0$ 的多步法
\[
  y_{n+1}-y_{n-1}=\frac{h}{3}(3f_n-f_{n-1}+4f_{n-2})
\]
是否收敛，并说明原因.

\medskip
\noindent\textbf{解}\quad
令 $n\to n+2$，将方法改写为标准线性 $k$ 步法形式：
\[
  y_{n+3}-y_{n+1}=\frac{h}{3}(3f_{n+2}-f_{n+1}+4f_n).
\]
这是 $k=3$ 步法，系数为
\[
  \alpha_3=1,\;\alpha_2=0,\;\alpha_1=-1,\;\alpha_0=0;\qquad
  \beta_3=0,\;\beta_2=1,\;\beta_1=-\frac{1}{3},\;\beta_0=\frac{4}{3}.
\]
对应的特征多项式为
\[
  \rho(\lambda)=\lambda^3-\lambda=\lambda(\lambda-1)(\lambda+1),\qquad
  \sigma(\lambda)=\lambda^2-\frac{1}{3}\lambda+\frac{4}{3}.
\]

\textbf{相容性检验.}\quad
\[
  \rho(1)=1-1=0\;\checkmark,\qquad
  \rho'(\lambda)=3\lambda^2-1,\quad\rho'(1)=2,\quad
  \sigma(1)=1-\frac{1}{3}+\frac{4}{3}=2\;\checkmark.
\]
相容条件 (5.6.8) 满足，故方法是\textbf{相容的}.

\textbf{零稳定性检验.}\quad
$\rho(\lambda)=\lambda(\lambda-1)(\lambda+1)$ 的根为
\[
  \lambda_1=0,\quad\lambda_2=1,\quad\lambda_3=-1.
\]
各根的模：$|\lambda_1|=0<1$，$|\lambda_2|=|\lambda_3|=1$.
位于单位圆周上的根 $\lambda_2=1$ 与 $\lambda_3=-1$ 均为\textbf{单重根}.

因此，$\rho(\lambda)$ 满足\textbf{特征根条件}，方法是\textbf{零稳定的}.

\textbf{结论.}\quad
方法既相容又零稳定，由定理 5.6.7，方法\textbf{收敛}.

\end{document}
